\chapter{Many-Valued and Paraconsistent Logic}
\label{ch:many-valued-paraconsistent}

\section{Logical matrices}
\label{sec:many-valued-paraconsistent-01}

\begin{definition}[Logical matrix]
A finite logical matrix for a propositional signature is the structure
\[
  \mathcal M=(\mathsf V,D,(f_c)_{c\in\Lang}),
\]
where \(\mathsf V\) is a finite nonempty set of values,
\(D\subseteq\mathsf V\) is the set of
designated values, and each \(n\)-ary connective \(c\) is interpreted by a
function \(f_c:\mathsf V^n\to\mathsf V\).  A valuation maps propositional
variables to \(\mathsf V\) and extends homomorphically to every formula.
\end{definition}

\begin{definition}[Matrix consequence]
For a formula set \(\Gamma\),
\[
  \Gamma\entails_{\mathcal M}A
\]
iff every valuation assigning a designated value to every member of
\(\Gamma\) also assigns a designated value to \(A\).
\end{definition}

The number of values does not determine a consequence relation.  The
operations and the designated set are equally essential.

